\documentclass[11pt]{article}
\usepackage[UTF8,fontset=fandol]{ctex}
\usepackage{microtype}
\usepackage[margin=1in]{geometry}
\usepackage{amsmath,amssymb}
\usepackage{graphicx}
\usepackage{booktabs,longtable}
\usepackage{array}
\usepackage{siunitx}
\usepackage[numbers,sort&compress]{natbib}
\usepackage[unicode,hidelinks]{hyperref}
\usepackage[capitalise,noabbrev]{cleveref}
\usepackage{xcolor}
\usepackage{url}
\graphicspath{{figures/}}
\setlength{\emergencystretch}{3em}
\setlength{\parindent}{2em}
\setlength{\parskip}{0.35em}
\providecommand{\tightlist}{\setlength{\itemsep}{0pt}\setlength{\parskip}{0pt}}
\renewcommand{\abstractname}{摘要}
\hypersetup{
  pdftitle={映射超表面全息术的技术管线:设计、制造与调制},
  pdfauthor={OptoMind}
}

\title{映射超表面全息术的技术管线:设计、制造与调制}
\author{OptoMind}
\date{}

\begin{document}
\maketitle

\begin{abstract}
基于超表面的全息术性能受三种耦合机制支配------纳米结构介导的相位控制、制造诱导的结构无序以及主动调制带宽------这些机制必须协同优化而非独立设计;要突破当前局限,必须缩小计算设计理想、可扩展制造现实与动态器件物理之间的差距。超表面全息术在静态波前工程方面展示了卓越的能力,但其向实时动态显示和先进成像系统的过渡尚不完整。本综述认为,该领域的核心瓶颈并非任何单一的技术组件,而是设计、制造和调制的解耦处理:逆向设计算法在优化理想化纳米结构时未考虑制造公差,纳米制造工艺引入了降低全息保真度的无序性,而主动调制机制则以牺牲衍射效率或光谱带宽为代价换取速度。通过围绕这三种耦合机制组织文献,本综述确定了独立优化成功之处、跨机制权衡主导之处,以及正在涌现以弥合这些差距的协同设计策略。综合结果表明,尽管单个组件已达到令人印象深刻的基准,但动态全息显示和集成成像中的系统级性能仍受限于效率、带宽、分辨率和刷新率之间未解决的张力------这些张力要求采用联合的设计-制造-调制方法,而非顺序优化。
\end{abstract}

\noindent\textbf{关键词：} 超表面全息术, 逆向设计超表面全息图, 机器学习超表面设计, 拓扑优化全息超表面, 纳米加工超表面全息图, 可扩展超表面制造, 动态超表面全息术, 可调谐超表面显示

\section{引言}\label{ux5f15ux8a00}

超表面全息术已成为纳米光子学领域的一种变革性范式,有望以能够任意调控波前的超薄平面器件取代庞大的光学系统。该技术的物理核心基于三种截然不同的相位调控机制:几何相位(Pancharatnam-Berry 相位)、传播相位(共振相位)以及混合方法。每种机制在效率、带宽、偏振敏感性和设计灵活性之间施加了根本不同的权衡,这些约束决定了所有下游应用的性能上限。虽然几何相位提供了稳健的偏振处理能力,但往往受限于色散限制;相反,传播相位虽能提供高效率,却因共振色散而受限于窄带宽。理解这些基础物理是驾驭现代全息系统复杂设计空间的先决条件。

光学神经网络(ODNN)研究正迅速演化为复杂的计算逆向设计框架,以克服基于直觉的设计方法的局限性。这些方法从基于梯度的拓扑优化 \citep{doi_10_2528_pier12020501_bb5007b1} 到先进的机器学习代理模型,旨在将所需的电磁响应直接映射到纳米结构几何形状上。最近的突破,如 MOCLIP 等基础模型,展示了在广阔设计空间中进行高通量、零样本预测的潜力,从而能够生成通过传统模拟难以处理的超表面图案 \citep{doi_10_1038_s41467_026_76714_x_8ca5d39b}。同样,包括扩散模型在内的生成式方法已被整合到工作流程中,以解决逆问题的非唯一性,为以低错误率初始化复杂几何形状提供了新途径 \citep{doi_10_48550_arxiv_2506_21748_dc49c062}。这一计算转变的重要性怎么强调都不为过:它有效地解锁了理论上无限的设计空间,使得实现自由曲面全息图和以前无法想象的多通道复用策略成为可能。

然而,在这些计算成功与物理现实之间仍存在一种关键的张力。大多数逆向设计算法建立在理想化模型之上,假设具有完美的几何控制能力,而忽略了纳米制造的随机现实。这种依赖造成了显著的''设计-现实差距'',即在计算机中生成的复杂图案在转化为物理器件时遇到了''制造上限''。电子束光刻等串行写入技术施加了严格的公差阈值,而纳米压印光刻等可扩展替代方案虽然在高体积制造方面前景广阔,但往往难以在大面积上保持纳米级保真度 \citep{doi_10_3390_encyclopedia5040197_fe6bfdc9}。此外,超表面向动态应用的扩展揭示了一个根本性的''速度-效率-带宽三元困境''。有源调制机制------涵盖相变材料、液晶和微机电系统(MEMS)执行器------能够实现实时重构,但不可避免地会引入衍射效率或光谱范围的损失 \citep{doi_10_1002_lpor_71186_5fa1be75}。目前没有任何单一平台能解决这一三元困境,迫使设计师接受特定的妥协,从而限制了动态全息显示的商业就绪程度。

本综述系统地绘制了基于超表面的全息术完整技术路线图,以解决这些脱节问题。我们追溯了从基本相位控制机制、计算逆向设计、纳米制造公差到有源调制的轨迹,并最终分析了成像和显示应用中的系统集成。我们的主要贡献是将解耦优化确定为核心瓶颈:即倾向于孤立地优化设计、制造和调制,而非通过统一框架进行协同优化的倾向。通过整合算法描述,并区分人类感知指标(用于显示)和机器性能指标(用于成像),我们阐明了每个应用领域面临的独特挑战。我们认为,未来的进步不取决于逆向设计算法或制造分辨率的孤立改进,而取决于整体协同设计策略,这些策略共同优化物理器件几何形状、制造约束以及下游计算重建过程。

本文的结构反映了这一技术路线图。第 S01 节建立了相位控制机制及其内在权衡的物理基础。第 S02 节详细阐述了计算逆向设计的演变,作为算法方法的权威参考,同时避免在后文中重复冗余内容。第 S03 节考察了这些设计的物理实现,重点关注制造引起的无序性和可扩展性限制。第 S04 节分析了有源调制机制及速度-效率-带宽三元困境。随后我们区分应用领域,其中第 S06 节专注于高保真成像和光学计算系统,接着第 S05 节处理动态显示架构和以人为中心的指标。最后,第 S07 节综合这些发现,提出跨领域的协同设计策略,明确指出在整个技术路线图中识别出的冗余和差距,以此证明需要集成优化框架。

\section{超表面全息术中的相位控制机制:基础与局限}\label{ux8d85ux8868ux9762ux5168ux606fux672fux4e2dux7684ux76f8ux4f4dux63a7ux5236ux673aux5236ux57faux7840ux4e0eux5c40ux9650}

超表面通过在纳米尺度上利用亚波长纳米结构单元调控光与物质的相互作用,实现了对光固有属性(包括相位和偏振)前所未有的控制。这种能力依赖于纳米制造技术的进步,如双光子聚合和聚焦离子束铣削,这些技术使得制造特征尺寸接近或低于电磁波长的复杂结构成为可能。这种波前工程的物理基础是自由空间中电磁场的传播,其由时谐场的亥姆霍兹方程描述;通过瑞利-索末菲衍射积分等方法求解该方程,可以描述次级球面子波如何干涉形成传播的波前。\citep{doi_10_21236_ada421453_63603a3f} \citep{doi_10_1364_ol_509688_f6e1ff85} 利用这些原理,超表面通过亚波长微/纳结构操控光场,实现了超越传统体块元件的轻量化、易集成和多维光学控制。然而,尽管这类材料提供的参数设计空间在实践上是无限的,但为满足多目标性能标准而探索这一空间仍然具有计算挑战性,这促使研究从基于直觉的参数扫描转向自动化逆向设计技术,以充分释放紧凑高效光学系统的潜力。\citep{doi_10_1021_acsphotonics_3c00572_fda8f742} \citep{doi_10_1002_advs_202504634_52bafb6c} \citep{doi_10_1038_s41377_026_02272_y_42e58df2}

几何相位(亦称 Pancharatnam-Berry 相位)代表了第一种基于亚波长光调控的波前工程机制。\citep{doi_10_1364_oe_27_015750_2ec65356} \citep{doi_10_1364_oe_568932_9e61cb80} 与源于谐振结构内光程变化的传播相位不同,几何相位依赖于各向异性纳米结构的空间旋转,从而产生与取向角两倍成正比的相移。该机制本质上将相位调控与偏振态耦合;具体而言,左旋和右旋圆偏振入射光会因其手性而产生不同的焦点或波前。代表性应用通过双焦点超透镜展示了这一原理:其中具有不同旋转角的单个单元结构可高效实现偏振依赖聚焦。\citep{doi_10_3390_s22051889_f6b907d2} 然而,由此产生的设计面临固有局限:由于相位响应取决于结构的衍射特性,此类器件往往仅在单一波长下有效工作,若无复杂的空间分割或多层堆叠,其在多光谱成像中的应用将受到限制。因此,尽管几何相位为偏振复用提供了一条直接途径,但其色散局限性 necessitates 引入下文将讨论的互补机制。\citep{doi_10_3390_s22051889_f6b907d2} \citep{doi_10_1117_1_apn_2_1_016010_57069347}

图 \ref{fig:visual_01}展示了本节讨论的代表性视觉证据。

\begin{figure}
\centering
\includegraphics[width=0.9\linewidth,height=\textheight,keepaspectratio,alt={Three-Panel Comparison of Metasurface Phase Control Mechanisms. Distinct phase mechanisms map to specific nanostructures and equations, creating unique efficiency-bandwidth trade-offs that define holographic performance limits. AI-assisted explanatory diagram; not empirical evidence.}]{figures/01_FIG-GEN-001.png}
\caption{Three-Panel Comparison of Metasurface Phase Control Mechanisms. Distinct phase mechanisms map to specific nanostructures and equations, creating unique efficiency-bandwidth trade-offs that define holographic performance limits. AI-assisted explanatory diagram; not empirical evidence.}\label{fig:visual_01}
\end{figure}

与源自坐标旋转的几何相位不同,传播相位依赖于单个超原子内的光程变化和共振延迟来积累所需的波前相移。该机制常利用连续域束缚态(BICs)或法诺(Fano)共振来实现高品质因子共振和强电场局域化,从而通过尖锐的色散光谱特征实现高效的相位调制。\citep{doi_10_1002_adpr_202400191_88f03cd9} 例如,混合金属-介电超表面可以通过抑制辐射损耗来支持这些BICs,其中对称性破缺将完美局域态转化为可调品质因子高达873的准BIC,用于主动传感应用。同样,全介电超表面利用高折射率捕获光学模式,并通过局域态与辐射连续体的干涉产生不对称的法诺轮廓,与金属对应物相比具有较低的耗散损耗。\citep{doi_10_1021_acsnano_4c16934_9c79f5cf} 然而,这种对共振光谱特征的依赖引入了一个基本约束:色散。由于相位响应本质上与波长相关的共振条件相关联,要在可见光谱范围内保持图像质量,通常需要复杂的缓解策略,例如为特定的RGB颜色分配单独的波导通道,而不是实现真正的单层消色差。这种由共振局域化提供的高效率与由色散限制的带宽之间的权衡,突显了后续混合方法旨在克服的物理边界。\citep{doi_10_1002_adpr_202400191_88f03cd9} \citep{s2_42f2fcc893f3a45ac750a8ffe67752f8533c0c61_58d9419c} \citep{doi_10_1038_s41377_025_01761_w_a73a6b84}

已发表的研究开发了混合金属 - 介质超表面,通过结合金属与介质组分,同时利用等离激元场局域和低损耗介质共振,旨在解决纯等离激元系统固有的效率损耗以及共振介质的色散挑战。该光学架构力求在强场增强与光学效率之间取得平衡,通常借助连续域束缚态(BIC)来抑制辐射损耗,同时保持高品质因子。例如,混合光子 - 等离激元非对称光栅已被证明能够通过几何结构和介电常数的对称性破缺,调控准连续域束缚态(quasi-BICs)与 BICs 之间的跃迁,通过抑制辐射损耗实现高达 873 的品质因子。在一种湿度调制策略中,由氧化铝和金膜上的交替介质光栅构成的混合光子 - 等离激元非对称光栅(HPAGs)利用聚乙烯醇来调控 quasi-BICs 与 BICs 之间的跃迁。随着相对湿度的增加,聚合物的体积膨胀和折射率降低会引发几何结构和介电常数的对称性破缺,从而抑制辐射损耗。\citep{doi_10_1002_adpr_202400191_88f03cd9} 作为对这些混合策略的补充,全介质超表面利用二氧化钛/石英或氮化镓/蓝宝石衬底等高折射率材料,为可见光波段的高光学效率提供了一条途径,相较于金属方案,其为相位控制提供了更坚实的物理基础。然而,即使是这些先进系统也面临局限;近红外波段的实际实施表明,残余相位失配和超原子透射率的根本限制可能会制约整体系统效率。这表明,虽然混合方案和高折射率介质方法缓解了此前的权衡取舍,但它们引入了与非局域行为和制造公差相关的新约束,必须在后续设计阶段加以应对。\citep{doi_10_1002_adpr_202400191_88f03cd9} \citep{doi_10_1063_5_0141881_93942631}

波前操控的根本能力最终取决于单个超原子的电磁特性,这些特性定义了系统可用的自由度,尽管混合策略和介电策略提供了平衡效率与约束的途径。然而,将期望的光学响应转化为物理几何结构的过程十分复杂,因为电磁参数与结构尺寸之间的映射存在奇异点突变,从而产生了非唯一性挑战,阻碍了传统设计方法的应用。近期的进展通过采用深度残差网络结合电磁特性分类,从透射系数预测超原子几何结构,有效缓解了由这些奇异性导致的不准确预测,并在数据集参数上实现了低于 0.23 的均方误差。\citep{doi_10_1109_iccem60619_2024_10559188_c3199408} \citep{doi_10_1109_isctis58954_2023_10213114_7fee3b45} 互补框架进一步利用表示学习将高维数据压缩为特征丰富的低维潜在空间,从而解决了设计与响应空间之间的多对一关系;在此空间中,耦合粒子群优化的变分自编码器能够以高达 94\% 的精度识别全局最优解。这些进展强调,相位控制机制的内在权衡需要复杂的逆向设计工具,以应对超原子工程中复杂且非唯一的景观,从而为从启发式方法向自动化计算框架的更广泛过渡奠定了基础。\citep{doi_10_1109_iccem60619_2024_10559188_c3199408} \citep{doi_10_2139_ssrn_4082312_28ca582d}

超表面工程已从传统的启发式方法根本性地转向自动化逆向设计框架,以应对元原子特性中固有的复杂权衡和非唯一性挑战。常规的模式到相位范式依赖于复杂、耗时的数值计算和专家直觉------这往往限制了可访问的自由度------而现代策略采用计算方法直接将所需电磁响应映射到最佳纳米结构几何形状。这一转变由集成多任务学习与遗传算法的深度神经网络架构所体现,对于氮化硅圆盘阵列,其正向预测准确率达到 96.7\%,逆向设计成功率达到 83.9\%,从而显著超越了传统神经网络的性能。同样,结合深度神经网络与遗传算法能够实现快速的相位到模式映射,在预设相位范围内控制种群生成,从而显著减少轨道角动量发生器和金属透镜等器件的优化迭代次数。此外,使用可调脉冲的新型时域逆向设计方法应运而生,以克服频域伴随优化相关的漫长仿真时间,促进了具有均匀宽带响应的高性能金属透镜的创建。通过这些数据驱动和算法工具取代繁琐的暴力搜索过程,研究人员现在可以访问更大的设计空间和更多样的电磁功能,直接解决前文概述的物理约束,并为学习非线性输入输出关系的具体深度学习实施奠定基础。\citep{doi_10_1016_j_rio_2026_100988_b00ee173} \citep{doi_10_1364_oe_478084_bf365dc5} \citep{doi_10_1109_apmc60911_2024_10867828_7c3f3b13}

深度学习模型现在解决了超表面逆向设计中映射非线性几何-响应关系这一核心挑战,这是在从启发式方法向自动化框架转变的基础上实现的。这些模型通过非线性激活函数和反向传播进行训练,能够从大规模数据集中智能地学习复杂的输入-输出相关性,从而在无需明确掌握支配物理规律的情况下,有效地逼近麦克斯韦方程组的解。这种数据驱动的方法通过用能快速预测单元特性的代理模型取代直觉性的微调,简化了双各向异性和多功能超表面的设计。高级实现利用潜在空间上的高斯混合模型或生成对抗网络来优化特定的子流形,从而实现将所需的全息相位分布精确映射到各向异性超原子的几何结构上。然而,由于需要生成大量任务特定的数据集(这些数据通常通过昂贵的时域有限差分法(FDTD)或严格耦合波分析(RCWA)等仿真获得),该框架的转变带来了显著的初始计算成本,且最终性能对数据集质量和超参数优化仍然敏感。最终,虽然深度学习扩展了超越传统边界的设计空间,但其实际部署需要在数据生成成本与复杂、高保真全息响应需求之间进行仔细平衡。\citep{doi_10_1515_nanoph_2019_0474_d2187501} \citep{doi_10_48550_arxiv_2204_00433_1c5cfd3e} \citep{doi_10_3390_photonics12100947_b44bdefd}

基于梯度的拓扑优化为探索模拟域内的材料分布以发现非直观的二元布局提供了一个严谨的框架,这与深度学习在映射非线性关系方面的优势相辅相成。该实现将设计空间视为离散网格,其中介电结构经过迭代更新以最大化特定的性能优值。其效率的核心在于伴随优化(adjoint optimization),这是一种通过每次迭代仅运行两次全波仿真(一次正向和一次伴随)来计算梯度的方法,无论设计参数的数量多少。为了确保这些连续的数学解能够转化为可行的制造结果,研究人员采用仿射几何变换和鲁棒性算法对轮廓进行二值化并过滤掉亚30纳米的特征。例如,此类方法已使热发射器通过超越传统谐振器的复杂自由曲面几何结构达到92\%的效率。此外,当与神经网络和傅里叶表示的损失函数结合使用时,这些逆向设计算法可以快速优化超表面几何结构以实现高保真的波前重建,即使在噪声条件下也能使均方误差接近10\^{}-6。然而,尽管这些技术有效地应对了前述效率与设计灵活性之间的权衡,它们仍必须面对逆散射问题固有的非唯一性挑战,这一问题将在下一节中讨论。\citep{doi_10_1002_adom_202500249_9967e710} \citep{doi_10_48550_arxiv_2603_15210_fdcfb92f} \citep{doi_10_1515_nanoph_2020_0579_8812c5b1}

基于梯度的方法和伴随方法为获得可行的制备形貌提供了高效途径,然而,它们所处的设计空间本质上仍存在不适定性,特别是在针对全息波前调控的超表面逆设计问题上。这一数学难题源于从期望的光谱响应反推物理结构时,必须解决相位与振幅耦合中固有的模糊性;具体表现为一种一对多映射,即单一目标光谱对应多种截然不同的几何构型。这种非唯一性会干扰确定性优化算法和传统的基于回归的神经网络,往往导致模式平均,并生成在通过麦克斯韦方程组求解器验证时无法产生目标光谱的无效几何结构。此外,实际的逆设计问题常常约束不足或动态演变,例如智能超表面面临不断变化的设计目标或更新的制备限制时。为应对这些挑战,表征学习提供了一种强有力的方法,它将搜索过程解耦为输入侧模型和输出侧模型:前者提供一个可复用且与物理机制无关的器件几何流形,后者则充当快速代理求解器。\citep{doi_10_2139_ssrn_4082312_28ca582d} 该策略使得能够在无需每一步评估都进行耗时的全波电磁仿真情况下,高效地探索候选方案并针对新目标进行快速重优化,从而在设计空间中稳定解的存在;否则,由于计算量呈指数级增长,全局优化方法在该空间中将变得计算上不可行。\citep{doi_10_71465_fapm546_bda368a5} \citep{doi_10_1002_adom_202502062_d73a85e1}

病态逆问题的求解为先进算法通过映射宏观复杂目标至微观几何结构以实际实现多通道全息术铺平了道路。深度学习辅助框架通过将设计空间分解来解决这一挑战;例如,特定方法采用分类器区分双自旋系统中的串扰类型,使人工神经网络能够拟合结构参数与相位响应之间的非线性函数,从而实现四通道涡旋复用。这种数据驱动的方法超越了简单的分类,延伸至全局优化领域,其中嵌入在物理传播模型中的进化算法搜索满足严格微分光场约束的最佳相位掩模。此外,电磁逆源方法已从二维推广至三维配置,利用垂直远场截面上的期望功率分布来推断等效面电流和宏观特性。通过将这些计算策略与混合相位机制提供的物理自由度相结合,研究人员可以克服散射问题的非唯一性,实现在自旋、波长和空间域上的精确波前调控,直接应对本章概述的基本权衡问题。\citep{doi_10_1117_1_apn_2_1_016010_57069347} \citep{doi_10_1117_12_3110717_26d5b63b} \citep{doi_10_1109_lawp_2023_3266337_3b44a88a}

\section{超表面全息图的计算逆设计:从前向模型到算法优化}\label{ux8d85ux8868ux9762ux5168ux606fux56feux7684ux8ba1ux7b97ux9006ux8bbeux8ba1ux4eceux524dux5411ux6a21ux578bux5230ux7b97ux6cd5ux4f18ux5316}

纳米光子器件设计已从基于直觉的经验启发式方法根本性地转向严格的计算逆设计,这一转变是由探索传统解析编码无法处理的高维参数空间的需求所驱动的。逆设计方法利用有效的数学方法解决优化问题,使其特别适合实现具有超越传统能力的卓越性能的复杂设计目标。这一演变既包括经典迭代技术(如拓扑优化,它以巨大的计算消耗换取高计算精度),也包括有望解决这些效率瓶颈的革命性深度学习方法。最早的优化方法代表了逆设计最古老且或许是最直观的应用,其中拓扑优化中的二进制表示和设计结构的重新参数化通过连续的性能优值最小化迭代,直至达到最优解。具体针对超表面全息图,ODNN研究已从简单的解析叠加发展到复杂的算法策略,如伴随梯度法,该方法通过求解复杂度与单次前向模拟相似的''伴随问题'',计算成本函数相对于所有设计变量的梯度。\citep{doi_10_48550_arxiv_2603_15210_fdcfb92f} \citep{doi_10_1038_s41598_024_67142_2_db612126} 尽管这些高级方法能够实现许多参数的并行优化并促进多层超原子等新型架构的发展,但它们通常运行在理想化的模型中,可能无法完全捕捉制造诱导无序的非凸性质或物理可实现性所需的约束。因此,虽然现代算法理论上解锁了复杂的波前控制,但其对完美几何假设的依赖造成了模拟最优性与现实制造器件保真度之间的关键差距。\citep{doi_10_1088_1361_6439_ad3a72_72d60f33} \citep{doi_10_1364_oe_27_030308_704a1c9b} \citep{s2_76a83ea3baed75ad88b7f0782eb9a694a330c853_cee16cac} \citep{doi_10_1515_nanoph_2021_0660_2d91a441}

逆设计从根本上逆转了传统工作流程,将目标的电磁性能参数直接映射到结构几何形状,标志着从经验直觉向严格优化的转变。前向设计通过简单的一对一关系从已知结构预测响应,而逆设计则面临一个数学上病态的挑战,即多个不同的几何配置可能产生相同的光谱。这种非唯一性造成了一对多的映射,混淆了确定性算法,并往往导致传统回归模型中出现模式平均或无效几何结构。为解决这一问题,现代框架采用了诸如串联神经网络(将逆预测器与固定前向模型耦合)或条件解码器(集成高斯噪声向量以在潜在空间中探索多样化解决方案)等策略。\citep{doi_10_1364_josab_497661_d333080a} \citep{doi_10_1038_s41467_025_57516_z_7c69e056} 此外,由于纯数据驱动的方法可能会生成看似有效但不遵循麦克斯韦方程组的结构,最近的进展引入了物理信息约束和光谱先验,以确保生成的设计在物理上是可实现的。尽管算法日益复杂,但对理想化映射的依赖突显了当前方法在面对现实制造中复杂非线性共振模式时的脆弱性,这为试图更稳健地导航这些连续参数空间的基于梯度的技术奠定了基础。\citep{doi_10_1088_1361_6463_ad3e05_d24e0a0a} \citep{doi_10_23919_piers_fall62445_2025_11394631_4fdb8517} \citep{doi_10_71465_fapm546_bda368a5}

基于梯度的拓扑优化作为一种迭代策略,用于探索逆问题公式中固有的广阔设计空间,旨在优化仿真域内的材料分布以满足特定的光学目标。该策略的核心是伴随方法(adjoint method),这是一种仅需两次全波仿真------一次正向和一次伴随------即可高效计算设计更新所需梯度信息的技巧,且其计算量与设计参数的数量无关。\citep{doi_10_1088_1367_2630_acfcd6_32444f17} 这种效率使得在具有任意色散材料的宽带范围内探索复杂、非直观的结构成为可能,远远超越了参数扫描或遗传算法等随机方法的能力,后者在高维空间中往往力不从心。例如,将此方法应用于热发射器时,已获得了效率高达 92\% 的自由形式二元布局,显著优于传统的谐振器设计。基于伴随方法的拓扑优化已被应用于亚衍射极限聚焦金属透镜的逆向设计,从而能够生成超越传统衍射极限的复杂结构。\citep{doi_10_1088_1367_2630_acfcd6_32444f17} 然而,在这些算法成功与物理实现之间出现了一个关键脱节:虽然优化过程通常作用于连续轮廓或中间密度值以促进平滑收敛,但最终输出必须经过二值化处理才能具备制造可行性。这一后处理步骤将连续变量强制转换为离散的材料状态,若反复迭代进行,可能会引入性能偏差并增加计算复杂性,这凸显了理想化的连续优化与纳米制造的离散约束之间的张力。\citep{doi_10_48550_arxiv_2603_15210_fdcfb92f} \citep{doi_10_1021_acsphotonics_3c00572_fda8f742} \citep{doi_10_1002_adom_202500249_9967e710}

深度学习与连续参数优化框架的结合通过采用全连接神经网络来近似正向散射函数,从而取代计算昂贵的电磁求解器,引发了一场范式转变。该设计依赖于代理模型,此处定义为一种经过训练的神经网络,能够从几何拓扑直接生成准确的电磁响应,而无需运行迭代仿真。\citep{doi_10_1109_jmmct_2025_3544270_6b4eb6dd} 通过学习这些映射关系,机器学习方法解决了传统设计方法的计算资源和时间限制,实现了按需的高效快速设计,其中结构直接从期望的目标计算得出。例如,基于物理的神经网络代理模型已与优化算法结合,用于搜索特定应用的相位库,如设计轨道角动量复用器,成功传输了高速信号。针对超表面的逆向设计,已开发出基于神经网络的代理模型,通过学习结构参数与光学响应之间的映射关系,为缓慢的电磁求解器提供了一种快速的替代方案。\citep{doi_10_1364_prj_450564_4b46735d} 虽然该策略有效规避了传统参数扫描所需的庞大电磁迭代成本,但它标志着从严格的物理求解向数据驱动近似的转变,为解析将光学响应映射回几何结构时固有的非唯一性所需的专用架构奠定了基础。\citep{doi_10_71465_fapm546_bda368a5} \citep{doi_10_1364_prj_450564_4b46735d} \citep{doi_10_1088_1361_6463_ab8036_d433363a}

标准代理模型加速了设计过程,但通常在解决将光学响应映射回几何结构时固有的非唯一性问题方面表现不佳,而先进的串联和级联架构则应对了这一挑战。串联网络通常将逆向预测器与前向模型配对以消除歧义,确保生成的结构产生预期的光学输出。延伸这一概念,级联神经网络架构由独立的前向和逆向设计网络组成,能够实现针对特定纳米光子功能的循环优化,有效解决了同时求解问题的复杂性。例如,某些具体实现利用深度残差网络结合电磁特性分类,从透射系数预测元原子几何结构,该方法减轻了映射过程中奇异点突变引起的误差,并将均方误差降低至小于 0.23。集成多任务学习与遗传算法的深度神经网络已被用于超表面的前向预测和逆向设计,允许从期望的透射系数准确确定元原子几何结构。\citep{doi_10_1016_j_rio_2026_100988_b00ee173} 同样,人们提出了循环框架,其中生成模型创建结构设计,而仿真神经网络评估其真实性,从而指导无严格结构限制的超表面发现。然而,这些复杂的方法并非没有局限性;级联方法可能仍会迫使收敛于多个有效解中的一个,且同步训练可能引入阻碍正确收敛的超参数依赖性。因此,尽管这些高级拓扑显著提高了全息生成的理论保真度,但它们仍在受限的解空间内运行,预示了接下来讨论的更广泛的泛化挑战。\citep{doi_10_1038_s41598_020_76400_y_aa62c15d} \citep{doi_10_1109_iccem60619_2024_10559188_c3199408}

尽管采用了串联和级联模型等先进网络架构来解决全息图生成中的非唯一性问题,一个根本性的约束依然存在:用于逆向设计的神经网络无法有效泛化到其训练数据集中特定信息分布之外的场景。这一局限性使得必须采用混合算法,将数据驱动方法与常规优化策略相结合,以弥补训练域外的性能差距。在此类混合框架中,神经网络执行粗略的全局搜索以估算期望解,而伴随方法等迭代优化算法则执行局部细化步骤以进一步提升性能,从而克服这两种技术各自的局限。为了具体解决在不从头重新训练的情况下将这些模型应用于不同功能的挑战,迁移学习已成为关键的赋能技术。通过微调特定的逆向设计网络参数同时冻结隐藏层,该方法实现了知识在异质超表面数据集之间的转移,例如从电磁诱导透明领域转向吸收或相位控制目标。\citep{doi_10_1364_ol_514212_108a4a7f} \citep{doi_10_1109_ieeeconf35879_2020_9330480_27bef614} 值得注意的是,所提出的网络仅需少至 700 个目标域样本即可完成训练过程,显著降低了跨域设计所需的数据阈值。这些混合与迁移学习策略代表了迈向稳健逆向设计的重要步骤,但它们仍局限于模拟环境中,往往忽略了制造过程中的随机现实,这种脱节最终限制了实际器件的保真度。\citep{doi_10_1021_acsphotonics_2c00968_06a2e96b} \citep{doi_10_1364_ol_514212_108a4a7f}

逆向设计方法现在积极利用超表面几何结构中的高自由度,以超越简单的相位控制来操纵光与物质的相互作用,这建立在缓解泛化限制的混合策略基础之上。该配置通过将复杂的结构参数视为可优化变量,超越了基于直觉的经验法则。例如,由多个厚度在几十到几百纳米之间堆叠层组成的光学多层薄膜结构,构成了一个组合优化问题,其中确定具有期望响应的最佳材料组合和厚度缺乏闭式解。现代算法通过自动识别最优结构来解决这一问题,而深度学习模型可以建立从光学响应到几何构型的直接映射,与传统优化相比,可在几分之一秒内完成设计。这些技术实现了特定的功能应用,例如使用逆向设计的介电超表面进行任意角度图像中继并进行像差校正,其中神经网络优化相位响应以补偿偏转图像中的严重畸变。此外,深度学习通过利用小尺度统计随机数据选择将单元统计特性映射到全阵列配置,从而促进了反射式超表面天线的逆向设计,避免了穷举参数扫描。将这些能力扩展到更大尺度,快速直接求解器现已实现厘米级圆柱形超表面的全波逆向设计,使得对因仿真成本高昂而此前难以处理的大面积全息元件进行计算优化成为可能。通过解锁这些复杂几何结构,计算工具将理论上的多功能性转化为实际的器件架构,尽管这种对精确几何定义的依赖为考察制造缺陷如何挑战这些理想化模型奠定了基础。\citep{doi_10_1016_j_isci_2025_112222_1f7c10f5} \citep{doi_10_1364_oe_519179_39046c11} \citep{doi_10_1002_mop_34068_b2e55b41} \citep{doi_10_1515_nanoph_2024_0127_c4fd5a12}

尽管逆向设计算法在利用高维自由度方面具备先进能力,但由于其依赖假设具有精确几何控制能力的理想化模型,因此出现了一个关键的设计与现实差距。优化策略通常依赖于对称性破缺微扰或连续域束缚态(quasi-BIC),其中光学寿命是在完美制造的假设下通过调整微扰强度来调节的。然而,启发式优化方法往往产生目标相位分布或传输矩阵,这些在物理上无法实现;当这些理想化设计被映射到可行的超表面结构时,这种不匹配会导致光学行为的显著偏差和性能下降。为缓解这一问题所做的努力,例如对相位分布应用平滑正则化以简化硬件设计,通常效果不佳,因为虽然更平滑的分布更容易实现,但它们会显著限制表达能力,且不能确保真正的物理可行性。同样地,用于光学超材料的物理信息机器学习模型往往仅最小化光学重建成本函数,而在损失函数中缺乏对设计参数的显式约束,这可能导致产生的解远离真实情况或实际可行性。虽然诸如循环内仿真训练等技术试图通过伴随方法来强制执行物理约束,但对于大规模系统而言,这些方法仍然昂贵得令人望而却步,使得许多优化后的设计容易受到制造诱导无序的影响。这种脱节强调了本章的核心论点:如果不将现实的制造约束直接整合到优化目标中,复杂的算法将继续产生理论上最优但实践中脆弱的器件。\citep{doi_10_1103_physrevapplied_21_044004_d1b93ffe} \citep{doi_10_48550_arxiv_2505_18377_a6fc9ec3} \citep{doi_10_1002_adpr_202300158_66d61883}

\section{超表面全息图的纳米制造:公差、无序与可扩展性}\label{ux8d85ux8868ux9762ux5168ux606fux56feux7684ux7eb3ux7c73ux5236ux9020ux516cux5deeux65e0ux5e8fux4e0eux53efux6269ux5c55ux6027}

包括深度学习和优化算法在内的逆向设计方法正日益被用于解决纳米光子学超表面全息图制造中固有的高维参数空间问题,有效地克服了传统基于直觉方法的局限性。\citep{doi_10_1109_metamaterials54993_2022_9920895_1a67346a} 这些方法生成的复杂连续或离散几何轮廓通常需要经过二值化和离散化步骤才能具备物理可制造性,从而立即依赖于电子束光刻等高精度纳米制造技术来实现所提出的结构。虽然经典迭代方法提供较高的计算精度,但在处理多目标或多尺度设计要求时面临巨大的计算消耗挑战,促使人们转向深度学习策略,后者有望在探索广阔的设计空间方面提供更高的效率。然而,从这些复杂的计算模型过渡到物理设备引入了一个关键瓶颈:由伴随梯度评估和神经网络产生的复杂几何形状所需的制造公差超出了当前制造能力的极限,这为分析目前为满足这些需求而采用的特定研究级工具奠定了基础。\citep{doi_10_1088_1361_6439_ad3a72_72d60f33} \citep{doi_10_48550_arxiv_2603_15210_fdcfb92f} \citep{doi_10_1109_metamaterials54993_2022_9920895_1a67346a}

当前的研究主要依赖高精度纳米加工技术,通过在纳米尺度调控光与物质的相互作用,来实现逆设计算法生成的复杂几何轮廓。具体而言,超表面全息图通常采用聚焦离子束刻蚀和各种光刻方法制备,这些技术能够制造特征尺寸接近或小于电磁波长的结构。例如,为了克服通常限制纳米结构高宽比和最小尺寸的纳米加工挑战,研究人员利用离子束刻蚀实现了工作在可见光波段的单片铌酸锂(LiNbO3)超表面。\citep{doi_10_1021_acsphotonics_1c00026_bdd15a92} 其中,电子束光刻(EBL)作为一种标准的实验室尺度工具,是一种串行写入技术,它利用聚焦电子束在基底上绘制图案,从而实现极高的分辨率。\citep{doi_10_1088_1361_6528_ad3c4c_ac8a5241} \citep{doi_10_1021_acsomega_4c01044_0ad8ee05} 尽管这些进展使得具有光束偏转和光场整形等新功能的紧凑光学系统成为可能,但确保高保真度的机制------即串行处理------也构成了关键瓶颈。由于这些方法是逐点或逐个特征地写入图案,而非并行处理,因此其吞吐量天生较低,仅能局限于小规模的演示验证。这种对串行制造的依赖造成了直接的矛盾:虽然计算方法可以探索几乎无限的设计空间以优化性能,但这些创新设计的物理实现却受制于当前制造工艺的缓慢速度,从而阻碍了从实验室概念向可扩展应用的转化。\citep{doi_10_1021_acsphotonics_3c00572_fda8f742} \citep{doi_10_1002_inf2_12116_b5c46838}

高分辨率串行写入技术能够实现复杂轮廓的初始原型制作,然而从计算逆向设计到物理超表面结构的转化不可避免地会引入结构无序性,本文将其定义为包括侧壁角度、特征粗糙度和位置误差在内的偏差,这些偏差会系统性地降低光学性能。这种退化在基于机器学习和优化的方法中尤为严重,因为在这些方法中,组成单元被赋予灵活性,可以采纳与邻居完全不同的基元以匹配所需的表面参数。尽管这种几何多样性对于实现特定的电磁响应是必要的,但它带来了显著的制造挑战;跨表面的物理结构突变会导致额外的互耦水平不可预测,而这些互耦并未包含在均匀化设计模型中。证据表明,虽然此类方法能够成功指导非均匀超表面的逆向设计,但由此产生的差异往往表现为远场辐射方向图中零值实现和旁瓣电平的误差。此外,采用伴随梯度评估和仿射几何变换的高级策略允许同时优化取向、位置和尺寸,然而,为了可行的制造而将这些连续轮廓二值化的要求增加了计算复杂性并可能导致保真度损失。因此,正是赋予现代逆向设计算法能力的几何自由度,在物理实现过程中成为了不确定性的来源,从而在理论预测与实验全息保真度之间形成了一道鸿沟。\citep{doi_10_48550_arxiv_2603_15210_fdcfb92f} \citep{doi_10_1109_tap_2021_3137496_d4a1d597}

除了物理实现过程中引入的结构无序之外,还存在一个更根本的障碍:算法本身所提出的设计方案往往本质上无法实现。用于光子超表面的逆向设计方法经常生成目标相位分布或传输矩阵,但由于严格的制备约束,这些目标在物理上无法构建。当将这些理想化的抽象模型映射到可行的超表面结构时,由此产生的失配会导致光学行为出现显著偏差并降低性能。虽然启发式方法试图通过施加平滑正则化来缓解这一问题以简化硬件设计,但此类措施通常不足;它们在限制系统表达能力的同时并未确保真正的物理可行性,往往仍会产生不切实际的设计方案。相反,闭环仿真训练通过利用伴随方法将优化直接嵌入训练循环中来强制保证可行性,然而由于该方法需要在每次迭代中求解麦克斯韦方程组,所得设计对于大型系统而言仍然成本过高且难以扩展。这种计算输出与物理现实之间持续的脱节凸显了一个关键差距:由于缺乏可扩展且具备物理感知能力的优化方案,算法潜力受到了抑制。\citep{doi_10_48550_arxiv_2505_18377_a6fc9ec3}

纳米压印光刻(NIL)等可扩展制造策略已被提出作为大规模生产的主要候选方案,旨在解决串行光刻方法固有的吞吐量瓶颈。NIL 基于直观的机械原理运作:将含有反向纳米图案的主模具压入抗蚀剂材料中,从而实现复杂超表面几何结构的并行复制,而非逐点写入。\citep{doi_10_3390_encyclopedia5040197_fe6bfdc9} \citep{doi_10_5772_intechopen_72534_9965e900} 该光学架构在理论上提供了一条途径,以克服与超表面集成特性相关的制造难度和成本的显著增加,这对于未来 3D 显示等需要超小型化和高质量光场调制的应用尤为重要。证据表明,虽然 TiO2 等材料可以提高工作效率,但其通过原子层沉积制备的过程耗时过长,难以实现商业化,这凸显了迫切需要更高效的微纳加工方法来解决这些生产问题。然而,尽管 NIL 和卷对卷加工有望实现更轻重量和更大视场的系统,但这些技术在米级面积上维持高效全息重构所需的纳米级保真度的能力仍 largely 未得到验证,导致提出的可扩展解决方案与最终商业化的严格要求之间存在关键差距。\citep{doi_10_1515_nanoph_2025_0269_b999ce93}

纳米压印光刻等可扩展路径在理论上实现了米级量产,然而现有的科学记录揭示了一个关键的验证缺失:证据严重偏向于计算设计方法,而缺乏对这些制造工艺在大规模应用中的实验验证。最近的进展展示了复杂的深度学习驱动框架,包括卷积神经网络和生成对抗网络,这些框架成功优化了超表面单元模式并捕获了全息成像的目标特征。例如,深度卷积生成对抗网络已被用于逆向设计电磁诱导透明超表面,证明了此类框架优化单元模式和捕获全息成像目标特征的能力。\citep{doi_10_1088_1402_4896_acf007_ad71688e} 同样,大规模光子逆向设计通过伴随优化和拓扑优化技术取得了进展,这些技术能够处理厘米级圆柱形超表面。该领域的进展包括使用多目标优化辅助的逆向设计方法进行超表面合成,说明了计算技术如何被应用于解决复杂的光子设计挑战。\citep{doi_10_1117_12_2529823_81c5e659} 然而,这些发展主要局限于计算领域,侧重于算法效率和潜在空间探索,而非工业复制所需的物理公差。因此,尽管能够理论上生成复杂自由曲面图案的设计工具已趋于成熟,但缺乏实验数据证实可扩展方法在从仿真过渡到米级制造时能否保持纳米级保真度。这种脱节强化了本章的中心论点,即超表面全息术的瓶颈已从设计能力转移到未经验证的可扩展制造现实,使得大面积器件的保真度成为一个悬而未决且关键的问题。\citep{doi_10_3390_photonics12100947_b44bdefd} \citep{doi_10_1515_nanoph_2024_0127_c4fd5a12}

\section{动态超表面全息图的主动调制机制:速度、效率与光谱权衡}\label{ux52a8ux6001ux8d85ux8868ux9762ux5168ux606fux56feux7684ux4e3bux52a8ux8c03ux5236ux673aux5236ux901fux5ea6ux6548ux7387ux4e0eux5149ux8c31ux6743ux8861}

传统超表面作为静态器件运行,意味着其光学能力在制造完成时即被永久固定。这种固有的''一次写入''特性造成了根本性的操作上限:尽管这些纳米结构表面能够提供强大的波前调控功能,却缺乏实时全息显示等动态场景所需的适应性。\citep{doi_10_1002_lpor_71186_5fa1be75} 尽管研究 efforts 曾试图通过集成活性材料或机械位移系统来克服这种刚性,但这些改进往往带来显著的代价,包括高光学损耗、制备复杂度增加,或依赖损害微型化外形的大体积致动器。因此,标准超表面无法在制备后重新配置其相位分布,这就 necessitates 探索专门的主动调制机制。解决这一局限性对于破解速度 - 效率 - 带宽三难困境至关重要,因为从静态操作向动态操作的转变构成了开发实用可重构全息系统的主要挑战。\citep{doi_10_1002_nap2_70185_56c8ab30}

相变材料(PCMs)通过结构相变改变折射率,为非易失性重构提供了一条途径,从而克服了传统超表面的静态局限性。\citep{doi_10_1515_nanoph_2024_0357_805aeaf6} 在动态非局域超表面中,该机制通过修改超原子参数(如旋转角)来实现多功能集成,无需重新设计整个晶格即可获得可调谐的光学响应。代表性证据展示了利用二氧化钒(VO2)实现的这一机制:当环境温度从 30 \(^\circ\) C 升高至 53 \(^\circ\) C 时,器件从绝缘态转变为能够进行窄带聚焦的动态金属透镜。例如,研究人员开发了一种 MoS2/VO2/Au/Si 超表面,利用 VO2 随温度变化的介电响应,动态调控光吸收率和结构色。对该器件加热会诱发绝缘体到金属的相变,改变 O-2p 态跃迁,从而实现对其光学性质的主动控制。\citep{doi_10_1515_nanoph_2023_0169_6ba026c6} 随着转换后的右旋圆偏振光围绕对应于这些温度变化的不同共振态进行动态聚焦,工作频率变得可调谐。然而,这种热调制带来了明显的代价:由于材料本征损耗增加,聚焦光的透射强度随温度升高而逐渐降低。因此,尽管 PCMs 提供了由单个超原子工程设计所决定的稳健光谱控制,但它们也引入了由吸收导致的可调性与衍射效率之间的权衡,这为可能缓解这些损耗机制的替代连续调谐方法奠定了基础。\citep{doi_10_1515_nanoph_2024_0357_805aeaf6}

图 \ref{fig:visual_02}展示了本节讨论的代表性视觉证据。

\begin{figure}
\centering
\includegraphics[width=0.9\linewidth,height=\textheight,keepaspectratio,alt={Active Modulation Platform Performance Comparison. No single active modulation platform currently satisfies the simultaneous speed, efficiency, and bandwidth requirements of the display-grade target zone. AI-assisted explanatory diagram; not empirical evidence.}]{figures/02_FIG-GEN-002.png}
\caption{Active Modulation Platform Performance Comparison. No single active modulation platform currently satisfies the simultaneous speed, efficiency, and bandwidth requirements of the display-grade target zone. AI-assisted explanatory diagram; not empirical evidence.}\label{fig:visual_02}
\end{figure}

混合超构光学系统通过将超表面与折射光学元件相结合以实现连续模拟控制,为高性能运行提供了一条替代路径,这与相变材料实现的非易失性开关形成了鲜明对比。在该架构中,传统折射元件提供主要光焦度,而集成超表面则作为校正器,用以抑制彗差和色差等像差,而这些像差通常是限制纯平面光学的瓶颈。这种协同机制实现了大空间 - 光谱带宽性能,有效解决了制约传统多透镜组件在长波红外热成像等应用中发展的权衡难题。代表性实施案例表明,交错排列超表面与折射组件可将平均聚焦效率从约 38\% 提升至超过 90\%,同时在宽光谱波段内实现超过 120 \(^\circ\) 的广视场。尽管这些混合设计显著增强了系统紧凑性并支持增强现实中的实时视觉任务,但它们也表明,克服带宽限制往往需要复杂的集成,而非依赖单一有源材料层,从而印证了本章的核心论点:为了平衡速度、效率和光谱范围,不同的架构折衷方案仍然是必要的。\citep{doi_10_1002_lpor_71431_d57da40c} \citep{doi_10_1364_ao_516890_6221deea}

超越液晶连续但光谱受限的调谐方式,弹性可重构超表面通过物理改变超原子几何结构以实现大幅相位调制,为动态全息术提供了一条独特的机械路径。这些系统基于一个直观原理:同时调控结构的有效质量和刚度,既能控制波的传播,又能保持阻抗匹配以实现高透射率。传统的依赖参数扫描或拓扑优化的设计方法,往往难以应对平衡这些多维变量所带来的计算负担,尤其是在追求多功能性能时。然而,最新进展表明,数据驱动的机器学习可作为有力工具,构建高维输入与输出之间所需的非线性映射,从而实现逆向设计,使元件在不牺牲透射效率的前提下覆盖完整的 2 \(\pi\) 相位范围。作为对上述通用机器学习策略的补充,空间耦合模理论结合可微分神经模拟器已成为一种特定框架,用于高效优化面向光谱成像和全息应用的超表面几何结构。该实现方案以高效近似替代传统电磁求解器,实现了光学编码器与重建网络的联合优化,相较于全波仿真,将设计周期缩短了一个数量级。尽管这种机械方法成功地将相位控制与电子方法中常见的吸收损耗解耦,但其本身也引入了物理致动所固有的速度限制,这进一步印证了本章的核心论点:实现动态可重构必然需要在调制深度、效率和响应时间之间进行权衡。\citep{doi_10_1016_j_matdes_2022_111560_e69e3350} \citep{doi_10_1021_acsphotonics_4c00171_f68cb79c} \citep{s2_76a83ea3baed75ad88b7f0782eb9a694a330c853_cee16cac}

为突破 MEMS 的机械限制,通过金属透镜进行的电子调制提供了一条利用亚波长微纳结构操控光场以实现超快开关的路径。这些器件提供了一个轻质且易于集成的平台,支持多维波前控制,从而推动了计算生成全息术和近眼显示领域的创新。热调制纳米光子相控阵(NPAs)已被开发为集成光源的全息显示器,与其他技术相比具有更高的刷新率。为解决热邻近效应(即加热一个像素会影响邻近像素的温度),邻近效应校正方法已扩展至仅相位优化的菲涅耳全息图。\citep{doi_10_1109_vr50410_2021_00058_601721f7} 与机械致动器不同,金属透镜依赖于静态结构几何或载流子动力学,而非物理运动,理论上允许快速重构。然而,尽管该策略在集成度和外形尺寸方面表现优异,却面临严格的光谱带宽限制。当前的制造框架限制了单个超原子所能实现的相位色散量,导致孔径尺寸、数值孔径和工作带宽之间存在相互制约。因此,对于大孔径器件而言,在整个连续光谱上实现严格消色差几乎是不可能的,迫使设计者接受非严格消色差的性能或减小孔径。这凸显了一个关键的效率 - 带宽权衡:虽然电子和结构调制解决了机械系统固有的速度瓶颈,但也引入了光谱范围和可扩展性方面的新限制,使其无法成为通用解决方案。\citep{doi_10_1038_s41377_026_02272_y_42e58df2}

超表面能够以前所未有的方式调控光的振幅、相位和偏振等属性,从而为 3D 显示器生成高质量的全息图和结构色。综合评估有源调制平台的性能揭示了一个根本性瓶颈:速度 - 效率 - 带宽三难困境。尽管这些平面光学器件在光学薄膜表面上设计空间非均匀性方面提供了高度的自由度,但要实现最终的商业化,仍需达到接近无限大的帧率;而目前的解决方案(如轨道角动量复用)仅能部分满足动态需求,使得这一目标仍遥不可及。证据表明,关键参数之间存在相互制约关系;具体而言,消色差超透镜的数值孔径、效率和工作带宽无法同时优化,且增加设计复杂度往往会引入串扰,从而降低工作效率。此外,材料选择也面临相互冲突的约束:低损耗材料(如 TiO2)虽能提高效率,但其制备过程耗时较长;而硅基替代方案则在衍射效率较低的情况下提供了不同的权衡。因此,目前尚无单一有源机制能在所有三个维度上占据主导地位,这证实了实用的动态全息显示必须在切换速度、光学效率和光谱范围之间应对固有的权衡,而非寄望于通用解决方案。\citep{doi_10_1515_nanoph_2025_0269_b999ce93} \citep{doi_10_1002_inf2_12116_b5c46838}

\section{超表面集成成像系统与光学计算:超越显示应用}\label{ux8d85ux8868ux9762ux96c6ux6210ux6210ux50cfux7cfbux7edfux4e0eux5149ux5b66ux8ba1ux7b97ux8d85ux8d8aux663eux793aux5e94ux7528}

逆向设计是一种直接从期望光学响应确定纳米结构几何形状的优化过程,它是推动超表面全息术从静态显示图案向动态计算成像转变的关键机制。\citep{doi_10_48550_arxiv_2506_21748_dc49c062} \citep{doi_10_1117_12_3099014_831c7349} 与模拟固定结构输出的传统正向建模不同,逆向设计能够探索高维参数空间,以实现超越传统能力的复杂多目标需求。传统方法采用拓扑优化或进化算法等迭代手段,虽然具有较高的计算精度,但在处理多尺度或多物理场挑战时,往往面临计算开销过大的问题。为克服这些效率瓶颈,数据驱动的机器学习方法作为一种革命性的替代方案应运而生;该方法利用深度神经网络构建代理模型,提供快速预测,从而规避了耗时的迭代电磁求解过程。这些框架通过将目标光学响应直接映射到纳米光子几何结构,实现了复杂超表面的高效构建,有效解决了设计复杂度日益增加与计算资源有限之间的矛盾。尽管经典方法在特定高精度任务中仍具价值,但集成深度学习对于实现先进非显示应用所需的复杂、像差补偿系统至关重要。\citep{doi_10_1088_1361_6439_ad3a72_72d60f33} \citep{doi_10_48550_arxiv_2204_00433_1c5cfd3e}

混合算法现在将神经网络与传统优化相结合,以改善超出特定训练分布的泛化能力,满足了逆向设计的基础需求。虽然纯数据驱动框架可以预测结构拓扑,但它们往往难以应对变化的材料属性或面外参数,因此需要一个将物理约束直接嵌入学习过程的统一方法。物理信息机器学习通过整合控制方程(如麦克斯韦波动方程)到损失函数或网络架构中作为正则化项来解决这一问题。\citep{doi_10_21070_acopen_10_2025_12939_5a404b4e} \citep{doi_10_1109_tmtt_2026_3658422_94a01e71} 该机制会对违反守恒定律或对称性要求的预测施加惩罚,有效地缩小了搜索空间,使模型能够在保持可解释性的同时,用较少的标注数据实现高精度。\citep{doi_10_71465_fapm546_bda368a5} 此外,由于逆向设计问题在数学上是病态的(因为存在一对多的映射,即不同的几何形状产生相同的光谱),物理信息扩散模型被引入以纳入光谱约束。这些模型确保生成的设计遵循基本的电磁定律,从而弥合仿真与现实之间的差距,否则几何上有效的结构可能无法产生目标光学响应。通过将物理先验知识与生成能力相结合,这些高级策略克服了黑盒深度学习的局限性,直接使得在实际金属透镜成像系统中校正像差所需的稳健设计成为可能。\citep{doi_10_1002_adom_202100548_6789fa4c} \citep{doi_10_1002_adpr_202300158_66d61883} \citep{doi_10_71465_fapm546_bda368a5}

从这些稳健的物理信息驱动设计策略转向实际部署,基于金属透镜的成像系统现在利用计算重建来主动补偿超表面固有的像差。这一转变使得无需运动部件即可实现高分辨率微型化显微镜和紧凑、稳定的三维层析成像,超越了仅靠被动光学元件所能达到的性能。虽然金属透镜通过定制化的超原子实现紧凑的光操控,具有独特优势,但由于基本制造限制,对于孔径超过一万波长的情况,通过相位匹配实现严格消色差几乎是不可能的。为解决这一问题,计算逆向设计框架利用角谱传播子评估输出强度和代价函数,从而优化级联超表面系统,避免了传统菲涅尔衍射积分中的近轴近似。\citep{doi_10_1364_oe_448757_33bdf93c} \citep{doi_10_1002_lpor_202401011_3926c17d} 该设计允许通过结构设计来减轻色差,而非仅仅依赖色散工程,促进了单片集成,从而产生极简光学系统,能够在扩展视场内实现消色差成像。此外,这类多元素级联配置(如三合镜)不仅限于简单的校正,还能表征复杂的三维光强分布,展示了此类系统在计算成像中进行高级波前整形的能力。通过将校正负担从纯粹的几何纳米结构设计转移到光学与算法的共同设计中,这些系统解决了孔径大小、数值孔径和工作带宽之间的权衡问题,直接推进了本章的核心论点:不同的应用领域需要耦合物理机制与算法重建的定制化策略。\citep{doi_10_3788_pi_2023_c01_9f71817a} \citep{doi_10_1364_oe_27_030308_704a1c9b} \citep{doi_10_1038_s41377_026_02272_y_42e58df2}

专用应用如今利用端到端协同设计,在标准超透镜成像的像差计算校正基础上,直接映射物理量而非仅仅形成视觉图像。在并行编码叠层成像中,一种集成了相位散射体和强度吸收体的无序工程化表面,解决了阻碍大规模阵列显微镜发展的扩展复杂性难题。这种即用型系统通过将工程化面永久附着于图像传感器,实现了比先前方案高出数个数量级的成像通量,从而消除了机械稳定性问题,并无需精密扫描即可实现开环光学采集。同样,超表面光学与非线性回归算法的端到端协同设计,通过将热辐射光谱直接映射为空间温度分布,实现了宽带温度成像。在 2024 研究中,研究人员展示了单层超表面前端与非线性

深度学习辅助的逆向设计现已能够独立控制自旋和频率等多个自由度,以实现复杂的多通道全息任务。由此产生的设计采用混合方法来解决串扰问题:首先由分类器区分特定的相位调制需求,然后让预训练的人工神经网络拟合在不同条件下物理参数与传播相位之间的非线性函数。例如,该框架促进了双自旋和双频率超表面的创建,这些超表面能够实现四通道离轴涡旋复用,其中高频参数首先确定,并作为低频参数自适应优化的预设条件。作为这种数据驱动策略的补充,非局域惠更斯超透镜提供了一种基于物理的高品质因子自旋复用成像途径,通过解耦相位和振幅响应,实现了对正交圆偏振光的独立操控,这对于结合明场和边缘增强成像且无光谱串扰至关重要。从学习复杂映射的生成对抗网络到潜在空间优化,这些进展共同证明了算法重建和定制化的逆向设计如何克服传统人工过程的局限性,直接支持了本章的核心论点:不同的应用领域需要耦合的物理和算法解决方案来管理分辨率和通道数量。\citep{doi_10_1117_1_apn_2_1_016010_57069347} \citep{doi_10_1038_s41377_024_01728_3_7dbdfab0} \citep{doi_10_3390_photonics12100947_b44bdefd}

超表面可被重构以通过光传播直接执行数学运算,这一范式被称为光学计算,其建立在为全息成像所确立的多通道控制基础之上。在此背景下,包括深度学习和耦合模理论在内的逆向设计技术,可优化介电结构以执行特定的信号处理任务,如傅里叶处理和矩阵乘法。例如,逆向设计技术已被用于基于介电超表面制备偏振不敏感的 C 波段达曼光栅,展示了优化后的介电结构执行特定信号处理任务的能力。\citep{doi_10_1016_j_rinp_2023_106238_0ebde0ff} 该设计过程将远场强度表述为几何参数的非线性函数,利用全连接神经网络拟合子函数,从而实现基于梯度的高效优化,以达到目标强度分布。针对谐振系统中局部优化往往失效的复杂性,先进策略采用结合了傅里叶表示拟合的复合损失函数,从而捕捉对于准确振荡行为至关重要的全局相位和振幅信息。这种计算严谨性延伸至离散散射体光学的大规模三维逆向设计,该方法能够管理实现任意光学功能(例如将光聚焦成离散的螺旋图案)所需的海量自由度。通过从成像向直接计算的转变,这些架构展示了定制的逆向设计如何将物理机制与算法目标相结合,从而克服传统正向设计的局限性。\citep{doi_10_1021_acsphotonics_4c00171_f68cb79c} \citep{doi_10_1515_nanoph_2020_0379_1cc596ab} \citep{doi_10_1117_12_2511323_41684364}

超越光学计算的矩阵运算,超表面在激光雷达(LiDAR)和传感领域的应用需要精确的波前准直和尖锐的光谱响应,以克服环境噪声。逆向设计的光子超结构通过优化光栅轮廓,直接产生面外宽束且均匀准直的光束,用于片上微流控传感,从而解决了传统光栅常因对称性导致的不均匀性等局限性。关于利用由准连续域束缚态(quasi-BIC)支配的法诺共振介电超表面进行成像的研究,阐明了这些逆向设计光子超结构如何解决传统光栅的局限性。通过优化光栅轮廓,此类系统旨在直接产生面外宽束且均匀准直的光束,这对于片上微流控传感等应用至关重要。\citep{s2_42f2fcc893f3a45ac750a8ffe67752f8533c0c61_58d9419c} 这种光场约束取代了笨重的台式激光器,降低了背景噪声,并实现了紧凑、高通量的生物分析。对于宽带偏振测量,探测能力依赖于对轴向色差的高效控制以减少串扰,这一挑战通过偏振依赖的相位优化得以解决,该优化允许在红外波长范围内进行直接、无需校准的测量,且不会出现早期消色差设计中常见的图像质量退化。此外,由准连续域束缚态(quasi-BIC)支配的共振介电超表面为高灵敏度任务提供了一种替代机制;通过扰动受对称性保护的状态以产生辐射通道,这些器件展现出不对称的法诺线型,与金属对应物相比,具有非凡的场增强效应和低耗散损耗。总之,这些不同的逆向设计策略表明,针对特定领域约束(无论是光束均匀性、光谱纯度还是共振质量)定制物理机制,对于推进计算成像系统至关重要。\citep{doi_10_1038_s41598_021_84841_2_2ce5ab57} \citep{doi_10_29026_oea_2022_220058_79ee78ee} \citep{s2_42f2fcc893f3a45ac750a8ffe67752f8533c0c61_58d9419c}

尽管逆设计方法已被证明能够实现波束偏转和谐振传感等特定功能,但将其扩展到多功能成像和计算所需的大规模设计空间仍是一个关键障碍。当前的方法往往难以在计算成本与全局最优性需求之间取得平衡,特别是在从受限的小参数子集转向现代纳米制造可实现的多样化形状和厚度时。为了解决这一可扩展性挑战,主动学习技术作为一种样本高效的策略应运而生,能够在保持多种光谱滤波优化任务性能的同时,将训练数据集的生成量减少高达 82\%。\citep{doi_10_1364_oe_559669_898a2205} 此外,基于深度学习的框架正通过混合策略演进以处理高维非凸问题;例如,将机器学习代理模型与凸优化相结合,可以在无需详尽全波仿真的情况下,实现满足复杂远场约束的多层非均匀超表面的确定性设计。进一步拓展该框架,像 MOCLIP 这样的基础模型正在开发中,以管理大规模的几何多样性,利用自由形式原型的随机生成和严格的实验验证来确保制造容差。然而,即使拥有这些先进的代理模型,在具有强层间耦合的系统中实现真正的全局最优性仍存在持续困难,因为相邻单元之间的不可预测相互作用可能会引入均质化模型未能捕捉到的差异。最终,从专用组件向稳健的计算成像系统过渡,取决于解决设计空间探索与物理可实现性之间的权衡。\citep{doi_10_1364_oe_559669_898a2205} \citep{doi_10_1038_s41467_026_76714_x_8ca5d39b} \citep{doi_10_1109_tap_2021_3137496_d4a1d597}

\section{超表面全息显示:架构、指标与应用准备情况}\label{ux8d85ux8868ux9762ux5168ux606fux663eux793aux67b6ux6784ux6307ux6807ux4e0eux5e94ux7528ux51c6ux5907ux60c5ux51b5}

在显示应用中,为实现复杂全息波前而应对高维参数空间,对传统解析方法而言仍具挑战。逆向设计方法通过利用有效的数学策略,将所需的光学响应直接映射到物理纳米结构几何形状上,从而实现了超越传统能力的目标,解决了这一计算需求。\citep{doi_10_1117_12_3099014_831c7349} \citep{doi_10_1364_ol_567547_28becc59} 为解决该映射中固有的非唯一性并弥合仿真与现实的差距,物理信息机器学习将麦克斯韦波动方程等控制物理方程,通过损失函数约束或残差建模直接整合到神经网络架构中。\citep{doi_10_1002_adpr_202300158_66d61883} 这种整合作为一种正则化手段,缩小了搜索空间,确保生成的结构遵守守恒定律,同时在较少标记数据的情况下提高了预测精度。尽管经典迭代方法提供较高的计算精度,但其巨大的计算消耗使其难以应对多目标超表面工程日益增长的复杂性,这使得深度学习驱动的框架成为从静态元件向动态系统原型演进的关键。\citep{doi_10_1088_1361_6439_ad3a72_72d60f33} \citep{doi_10_1002_adpr_202300158_66d61883} \citep{doi_10_71465_fapm546_bda368a5}

基础的逆设计方法将所需的光学响应映射到结构参数,然而标准神经网络往往无法在其训练数据集之外泛化,限制了其在新型全息轮廓中的实用性。为了克服这一局限,混合算法应运而生,将数据驱动模型与严格的优化技术相结合,第 2 节详细阐述了该策略在串扰和刷新率等显示约束方面的具体应用。通过利用这些先进的计算框架来导航多物理场耦合和复杂的设计空间,研究人员可以解决波前工程中固有的非唯一性和效率瓶颈。\citep{doi_10_1038_s41467_026_76714_x_8ca5d39b} 因此,这些算法的演进解决了核心计算挑战,为下一个关键障碍铺平了道路:将这些优化轮廓转化为具体的透射式、反射式或混合系统架构。\citep{doi_10_1021_acsphotonics_2c00968_06a2e96b} \citep{doi_10_1038_s41467_026_76714_x_8ca5d39b} \citep{doi_10_5757_asct_2023_32_6_141_32febf1c}

超表面全息显示主要通过透射式或反射式架构实现,这些架构作为核心波前调制元件,集成于完整的静态和动态投影系统中。尽管这些独立配置利用超小型化实现了更轻的重量和更大的视场,但它们往往在数值孔径、效率和工作带宽之间面临相互制约,直接损害了亮度均匀性和色域等人类感知指标。为解决这些局限性,混合超表面 - 折射透镜架构已成为一种可行的集成途径,可针对显示需求实现紧凑、高质量的成像。在该光学架构中,折射元件提供主要光焦度,而超表面则被策略性地交错排布,以校正对视觉保真度至关重要的光学像差。这种分工得益于专用的逆向设计算法,例如标量场射线 - 波混合传播方法,该方法能够在传统像空间之外优化纳米结构。例如,在中红外波段透镜中,将超表面同时置于物空间和像空间已被证明可将平均聚焦效率从 38\% 提升至超过 90\%,显著优于仅将超表面置于像空间的设计。尽管在架构效率方面取得了这些进展,但向最终商业化的过渡仍受到制造复杂性以及实现连续 3D 全息术(而非简单的带深度线索的 2D 投影)这一固有挑战的阻碍。\citep{doi_10_1515_nanoph_2025_0269_b999ce93} \citep{doi_10_1364_ao_516890_6221deea}

利用上一节建立的混合架构,基于深度学习的逆向设计现在能够将所需的电磁响应精确映射到特定的结构参数,从而克服了传统解析方法的局限性。该方法通过采用结合正向和逆向预测网络的循环结构,解决了输入与输出之间固有的''一对多''关系,与传统流程相比,不仅提高了精度,还显著节省了设计时间。具体而言,这些模型已成功应用于合成用于线偏振至圆偏振控制的定制化单层散射体,有效解决了复杂的非线性映射问题。深度学习已被用于逆向设计透射型超表面,以实现线偏振至圆偏振控制,有效解决了此类转换所需的复杂非线性映射。\citep{doi_10_1088_1361_6463_acefdf_20072f6c} 将这种能力扩展至多通道系统,深度学习辅助框架促进了双自旋和双频率超表面的构建,这些超表面能够通过涡旋复用产生四通道离轴全息输出。深度学习辅助框架已使得构建能够产生四通道离轴全息输出的双自旋和双频率超表面成为可能。\citep{doi_10_1117_1_apn_2_1_016010_57069347} 通过利用分类器区分相位调制需求,并利用人工神经网络拟合不同串扰条件下的非线性函数,这些设计可自动确定复杂波前整形所需的物理参数。虽然这些进展展示了通往多波长和多偏振显示架构的清晰路径,但它们也引入了关于串扰和制造模糊性的系统级权衡,这将在下一部分进行探讨。\citep{doi_10_1088_1361_6463_acefdf_20072f6c} \citep{doi_10_1117_1_apn_2_1_016010_57069347}

功能复用能够实现复杂的全息输出,然而驱动这些架构的逆向设计过程面临一个根本性瓶颈:期望的电磁响应与物理结构参数之间映射的非唯一性。这种一对多的关系使确定性制造变得复杂,因为不同的超原子几何结构可能产生几乎相同的光学特性,从而在选择可制造方案时造成歧义。证据表明,若缺乏严格的数据集筛选,在此类数据上训练的神经网络将难以收敛,这是因为对应不同结构变量的相似电磁特征会在训练过程中引入不稳定性。强制建立一一映射的策略通常涉及计算欧几里得距离以剔除冗余数据点,但若阈值设定过于激进,该实施方法存在丢弃有效设计的风险。此外,随着设计复杂度增加以支持更高的自由度,通道间的串扰成为一个关键障碍,不仅降低了工作效率,还 necessitating 与传统光学的混合集成,而这可能会推高生产成本。这些挑战凸显出,主要约束不再仅仅是实现特定的波前,而是确保逆向设计能生成一种独特、可制造的结构,且该结构能够在无通道间干扰的情况下运行。\citep{doi_10_1088_1361_6463_ad3e05_d24e0a0a} \citep{doi_10_1515_nanoph_2025_0269_b999ce93} \citep{doi_10_1038_s41467_026_76714_x_8ca5d39b}

突破此前识别的静态设计瓶颈,动态超表面系统现利用对称性破缺机制实现实时自适应。一种关键方法涉及混合金属-介电超表面,其支持连续域束缚态(BIC),这是一种共振态,由于辐射损耗被抑制而具有高品质因子和强电场局域化特性。尽管在准连续域束缚态(quasi-BIC)与真实BIC之间转换仍具挑战性,但近期研究表明,打破几何和介电常数对称性可以动态调控这一转变。具体而言,湿度调制策略利用吸湿材料(如聚乙烯醇),随着相对湿度的增加,这些材料发生体积膨胀并降低折射率,从而改变非对称光栅的面外对称性,将共振品质因子调节至 873。\citep{doi_10_1002_adpr_202400191_88f03cd9} \citep{doi_10_1364_oe_574568_72469f3c} 作为对这种固态调制的补充,基于微流体的架构为宽带动态控制提供了另一条独特路径。通过将电磁设计与层流两相流分析相结合,实验报告实现了自适应迷彩斗篷,其中镓基液态金属在微通道内被氢氧化钠溶液置换。基于微流体的实验工作展示了宽带动态自适应迷彩斗篷,该系统依赖多物理场耦合分析来管理微通道内镓基液态金属被氢氧化钠溶液置换的过程。\citep{doi_10_1002_lpor_202500327_a3bf6deb} 这种多物理场集成实现了可重构散射图案和多幻象操作,证明流体驱动可以有效修改近场分布和远场响应。综上所述,这些对称性破缺和微流体方法说明了向有源组件的转变,但也表明,要实现商业显示所需的刷新率和眼盒均匀性,仍需进一步进行系统级集成。\citep{doi_10_1002_adpr_202400191_88f03cd9} \citep{doi_10_1002_lpor_202500327_a3bf6deb}

近期消色差超表面波导在增强现实(AR)领域的进展,最能体现从动态调制机制向商业化应用就绪的过渡。这些系统利用集成在单板波导上的超薄氮化硅(SiN)二元纳米结构耦合器,为全彩显示提供了一条简单且具成本效益的制备路径。针对 AR 显示的消色差超表面波导研究强调,采用集成在单板波导上的超薄氮化硅二元纳米结构耦合器,是实现全彩应用的简便且低成本的制造方案。\citep{doi_10_1038_s41377_025_01761_w_a73a6b84} 该机制依赖于优化光学响应,以在反射模式下有效利用红、绿、蓝波长的不同衍射级,从而在解决衍射波导中长期存在的色差问题的同时,实现高色彩和角度均匀性。然而,尽管在紧凑外形和大视场支持方面展现出独特优势,当前原型机仍面临特定的系统级限制,阻碍了其即时部署。所用 RGB 光与设计波长之间的失配,以及原型机固有的像差,限制了这些设计潜力的充分发挥。此外,这些系统固有的偏振选择性要求未来开发二维偏振不敏感超表面,以拓宽其适用范围。要实现更高的效率,还需要对输入耦合器进行战略优化(例如通过反射涂层),同时抑制过度高阶衍射对于实现宽视场显示至关重要。因此,虽然逆向设计促成了这些功能架构的实现,但 ODNN 研究的主要瓶颈已从单组件性能转向系统级集成,即需要同时解决效率、偏振无关性和衍射控制问题,以实现实时、全色域的 AR 体验。\citep{doi_10_1038_s41377_025_01761_w_a73a6b84}

\section{超表面全息术中的跨领域挑战与协同设计策略:解决耦合机制瓶颈}\label{ux8d85ux8868ux9762ux5168ux606fux672fux4e2dux7684ux8de8ux9886ux57dfux6311ux6218ux4e0eux534fux540cux8bbeux8ba1ux7b56ux7565ux89e3ux51b3ux8026ux5408ux673aux5236ux74f6ux9888}

逆向设计方法利用有效的数学手段来解决优化问题,使其特别适用于探索高维设计空间,并实现通过基于直觉的参数化无法触及的高复杂度设计目标。虽然经典的迭代方法具有较高的计算精度,但其计算消耗过大,从而催生了作为快速前向预测和图案到相位映射革命性工具的深度学习框架。例如,特定的深度学习架构已展现出 96.7\% 的前向预测准确率和 83.9\% 的逆向设计成功率,在效率上显著超越了传统神经网络。然而,超表面工程的核心挑战在于,由于非唯一性,该逆问题在数学上是病态的,即多种不同的几何构型可能产生相同的光谱。\citep{doi_10_30560_ijas_v9n1p1_31843ad0} \citep{doi_10_1364_oe_419138_33b8279d} 这种一对多的映射使确定性优化算法和基于回归的传统神经网络陷入困境,往往导致模式平均化,并在经麦克斯韦求解器验证时生成无法产生目标光谱的无效几何结构。因此,尽管数据驱动方法提供了速度,但它们缺乏对底层物理机制的内在理解,需要引入稀疏性假设或物理信息约束来解决这些根本性的歧义,从而实现计算上可行的解决方案。\citep{doi_10_1088_1361_6439_ad3a72_72d60f33} \citep{doi_10_1016_j_rio_2026_100988_b00ee173} \citep{doi_10_1515_nanoph_2019_0474_d2187501} \citep{doi_10_71465_fapm546_bda368a5}

深度学习为导航广阔的设计空间提供了强大工具,但纯粹的数据驱动或启发式方法往往因忽视波动物理的基本约束而无法产生物理上可实现的超表面。构建转移矩阵的启发式方法通常过度简化复杂的光与物质相互作用,导致理想化的相位分布无法在实际几何结构中实现,并在投影到可行结构时造成显著的性能下降。这一局限性源于缺乏显式设计参数约束的逆模型仅学习最小化光学重建误差,实质上是将问题视为忽略可制造性和物理有效性的黑盒映射。为解决此问题,物理信息机器学习通过将麦克斯韦方程等领域知识直接整合到学习过程中(通过损失函数、残差建模或架构初始化)来发挥作用。\citep{doi_10_1007_s40314_026_03818_x_05a7eee9} \citep{doi_10_48550_arxiv_2606_04143_0a5e8b07} 通过将这些支配性的物理原理嵌入为正则项,模型被引导至一个本质上满足守恒律和对称性要求的解子空间,从而提高泛化能力并减少对大规模标注数据集的需求。这一策略与可能收敛于非物理局部最优的无约束模型形成鲜明对比,确保最终设计在光学上高效且稳健地符合实际制造所需的底层物理规律。\citep{doi_10_1002_adom_202502062_d73a85e1} \citep{doi_10_1002_adpr_202300158_66d61883} \citep{doi_10_48550_arxiv_2505_18377_a6fc9ec3}

图 \ref{fig:visual_03}展示了本节讨论的代表性视觉证据。

\begin{figure}
\centering
\includegraphics[width=0.9\linewidth,height=\textheight,keepaspectratio,alt={Coupled Mechanism Framework for Metasurface Holography. System performance boundaries emerge from synergistic couplings between phase selection, fabrication disorder, and modulation constraints, necessitating co-design interventions. AI-assisted explanatory diagram; not empirical evidence.}]{figures/03_FIG-GEN-003.png}
\caption{Coupled Mechanism Framework for Metasurface Holography. System performance boundaries emerge from synergistic couplings between phase selection, fabrication disorder, and modulation constraints, necessitating co-design interventions. AI-assisted explanatory diagram; not empirical evidence.}\label{fig:visual_03}
\end{figure}

为解决超表面逆向设计中固有的非唯一性与收敛瓶颈,先进的深度学习架构如今提供了本节核心的协同设计框架所需的具体机制。级联神经网络通过训练一个逆向网络以最小化重预测光谱的误差(而非直接的几何差异),从而解决了病态映射问题,有效规避了不同几何结构产生相同光学响应这一''一对多''难题。\citep{doi_10_1063_5_0171437_92a49ec0} \citep{doi_10_1515_nanoph_2021_0434_1357d09e} 生成式方法已从常受模式坍塌困扰的生成对抗网络(GANs),进一步演变为物理信息扩散模型;后者通过逆转渐进噪声过程,在严格的光谱约束下生成多样化的超原子形状。\citep{doi_10_71465_fapm546_bda368a5} 针对电磁特性映射中的奇异点突变问题,结合特性分类的深度残差网络已展现出预测几何结构的能力,其均方误差低于 0.23,从而稳定了全息成像超表面的设计。在基础模型尺度上,MOCLIP 等框架利用超过 466,000 组几何 - 光谱对的庞大实验数据集学习共享潜在空间,实现了约 200,000 光谱/秒的零样本预测速度,同时纳入了纯模拟数据所缺失的制造容差。作为对这些数据驱动策略的补充,伴随变量法(AVM)提供了一种基于梯度的优化框架,能够高效获取复杂几何结构的灵敏度信息,加速最优设计的搜索过程,而无需完全依赖大型训练语料库。\citep{doi_10_2528_pier12020501_bb5007b1} \citep{doi_10_1109_nemo62710_2025_11215427_067e96bb} 综上所述,这些架构将范式从缓慢的启发式微调转变为几何结构与光谱响应的快速联合优化,通过确保计算速度不以牺牲物理有效性或可制造性为代价,直接解决了耦合机制瓶颈。\citep{doi_10_1038_s41467_026_76714_x_8ca5d39b} \citep{doi_10_71465_fapm546_bda368a5} \citep{doi_10_1109_iccem60619_2024_10559188_c3199408} \citep{doi_10_1109_pn62551_2024_10621761_0bc70b78}

基本的优化范式决定了在高维非凸设计空间中导航时,收敛速度与全局最优性之间的平衡。逆向设计方法目前正分化为基于梯度的方法(如伴随拓扑优化)和无梯度的自然启发算法。基于梯度的方法利用互易定理计算性能优值相对于设计参数的变化率,从而允许同时更新整个设计向量以快速达到局部最优;然而,这种策略本质上具有局部搜索特性,存在将解陷入次优极小值的風險。相反,无梯度技术通常受狼群狩猎等自然现象启发,整合了不同的探索与开发阶段,无需显式梯度信息即可跳出这些局部陷阱,但其计算成本往往较高。传统的启发式方法(如遗传算法和粒子群优化)进一步说明了这一瓶颈,因为它们计算量大且不可微,与现代可微代理模型相比,阻碍了其与下游任务的全流程集成。这种二分法表明,虽然基于梯度的方法为定制光学功能提供了效率,但在复杂的相位景观中难以保证全局最优性;而无梯度方法则以速度为代价提供了更广泛的搜索能力,凸显了解决这些耦合机制瓶颈的策略需求。\citep{s2_2516273dc5b4511c838bf9236a61f71f4e32f6b5_1724698f} \citep{doi_10_1088_1367_2630_acfcd6_32444f17} \citep{s2_76a83ea3baed75ad88b7f0782eb9a694a330c853_cee16cac}

在高维空间中导航的优化范式遇到了一个基本的''耦合机制瓶颈'',其中几何相位和传播相位通过共享的结构参数紧密相连。例如,在双自旋和频率复用超表面中,某一频率下的相位响应可能取决于纳米光栅尺寸和取向角的一个特定子集,而另一频率下的响应则由这组相同变量的全集决定。\citep{doi_10_1117_1_apn_2_1_016010_57069347} \citep{doi_10_1002_advs_202504634_52bafb6c} 这种相互依赖性创造了一个复杂的优化景观,调整一个参数以满足一个通道会无意中改变另一个通道的响应,从而引入显著的串扰。当前的逆向设计方法通常试图通过将超表面视为孤立组件来解决这一问题,利用分类器和迭代神经网络根据预设条件抑制特定的相位串扰。然而,此类方法往往未能考虑更广泛的全局系统级相互作用(如相邻元原子之间的耦合或基底效应),因为它们依赖于拟合目标振荡行为的损失函数,而未充分嵌入整个光学环境的物理约束。因此,孤立地处理设备限制了仿真到实验工作流的预测准确性,加强了本章的核心论点:解决这些权衡需要从组件级优化转向集成协同设计框架。\citep{doi_10_1117_1_apn_2_1_016010_57069347} \citep{doi_10_1515_nanoph_2020_0379_1cc596ab}

协同设计框架已成为解决前述耦合机制瓶颈的主要策略,此处定义为将计算逆向设计与制造感知约束直接整合的联合优化协议。其直观前提是,高保真相位控制无法脱离制造无序的物理现实而实现;相反,这些容差必须嵌入到优化循环本身之中。从机制上讲,这是通过构建联合损失函数或统一的机器学习架构来实现的,这些架构考虑了谐振组件固有的非唯一性和振荡行为。例如,最近的方法利用分类器区分几何相位和传播相位之间的特定串扰条件,从而实现对元原子参数的自适应优化,以满足多通道波前整形要求。同样,统一框架通过条件解码器映射特征向量,并采用预训练的前向模型来减轻通常会恶化收敛性的非唯一性效应,从而解决设计规格不完整的问题。这些方法显示出显著的效率提升,与传统遗传算法相比缩短了设计时间,同时在扩展的参数范围内保持了鲁棒性。然而,尽管此类策略有效地弥合了理想化仿真与实验可行性之间的差距,但对于最佳联合优化协议的共识仍不完整,突显了接下来讨论的更广泛系统集成之必要性。\citep{doi_10_1117_1_apn_2_1_016010_57069347} \citep{doi_10_1515_nanoph_2020_0379_1cc596ab} \citep{doi_10_23919_piers_fall62445_2025_11394631_4fdb8517}

协同设计必须进一步扩展到端到端系统优化,其中超表面几何结构与下游计算算法被联合调整以克服物理限制。在此框架下,光学元件无需独立形成完美图像;相反,它以特定后处理算法能够最优重构的方式编码信息。例如,最近的工作表明,将单层超表面与非线性普朗克回归算法进行协同优化,能够从长波红外辐射中高质量地重构温度图,从而有效绕过通常要求笨重多层光学的延迟-带宽约束。这种联合设计依赖于可微电磁仿真器,它将几何参数直接链接到复反射和透射系数,并在优化过程中准确考虑薄膜层内的多次内部反射。此外,混合超光学系统将这些平面元件与传统折射透镜相结合,以利用两个平台的优势:折射元件提供主体光学功率,而超表面校正像差,从而实现适合增强现实热成像的大空间-光谱带宽。通过将传感器、光学系统和算法视为一个统一的可微管道,这些策略解决了困扰独立优化子系统的微型化与性能之间的权衡问题,直接回应了本章的核心论点,即打破耦合机制瓶颈需要整体设计框架。\citep{doi_10_1002_adom_202402498_4083d234} \citep{doi_10_1002_lpor_71431_d57da40c}

新兴物理原理为通过高品质因子共振从根本上突破性能极限提供了途径,这与解决架构权衡的系统级协同设计相辅相成。对称性保护的连续域束缚态(BIC)代表理论上与自由空间解耦的模式;通过引入受控的对称性破缺微扰(如二聚化),这些模式演变为具有可调品质因子(Q 因子)的泄漏准连续域束缚态(quasi-BIC),其 Q 因子与微扰强度成反比缩放。该机制在保持强场局域的同时实现了精确的波前控制,正如环形偶极子驱动的超表面所证实的那样:此类器件实现了完全的偏振选择性,并在大角度斜入射下稳定运行,优于以往仅限于正入射的设计。环形偶极子驱动的准连续域束缚态已被用于介电超表面中的大角度近场图像显示,实现了完全的偏振选择性和斜入射下的稳定运行。\citep{s2_5a46bf1115a13c8e88d03d52a56e80871e3550cd_094f009d} 与这些物理进展相辅相成的是,基于欧拉拉格朗日潜行动力学的频率传递技术这一计算前沿,提出了在复杂多物理场耦合约束下协同优化超表面的方法。综上所述,这些进展表明,解决耦合机制瓶颈不仅需要整合制造容差,还需利用奇异共振态和先进的逆向设计算法,以联合优化光学效率、角度鲁棒性和系统功能。\citep{doi_10_1103_physrevapplied_21_044004_d1b93ffe} \citep{s2_5a46bf1115a13c8e88d03d52a56e80871e3550cd_094f009d} \citep{doi_10_1038_s41467_025_57516_z_7c69e056}

\section{结论}\label{ux7ed3ux8bba}

本综述系统地绘制了基于超表面的全息术的技术管线,追踪了其从基本相位控制机制到复杂系统级应用的发展轨迹。分析得出的核心发现是,虽然计算逆向设计已有效地解锁了理论上无限的设计空间,但高保真度动态全息系统的实际实现仍受到持续的''设计-现实差距''的制约。这一差距属于增量工程改进的范畴,源于当前开发生命周期中对物理机制、制造公差和系统集成进行的未解耦处理。

超表面全息术依赖于三种不同的相位控制机制------几何相位、传播相位以及混合方法,每种机制都在效率、带宽和偏振敏感性之间施加了内在的权衡。如前所述,几何相位提供了稳健的偏振控制,但存在色散限制;而传播相位则以窄带宽为代价提供高效率,这是由于谐振色散所致。尽管混合策略试图在这些物理边界中导航,但它们引入了设计非唯一性的新复杂性。为了解决这些挑战,光学深度学习(ODNN)研究已从基于直觉的经验法则转向复杂的计算逆向设计框架,包括深度学习代理模型和基于伴随变量的拓扑优化。这些算法成功地将所需的电磁响应映射到纳米结构几何形状上,克服了逆散射问题的不适定性。然而,这些计算上的成功往往依赖于假设完美几何控制的理想化模型,忽略了纳米制造的随机现实。

这种对理想化的依赖在向物理实现过渡时造成了关键瓶颈。综述指出,当前的纳米制造工艺,特别是电子束光刻等串行写入技术,施加了严格的公差阈值,导致全息保真度系统性地低于理论预测值。虽然纳米压印光刻等可扩展方法为大规模生产提供了一条路径,但在维持米级面积上的纳米级保真度方面,实验验证存在显著不足。因此,逆向设计算法生成的复杂自由曲面图案经常遭遇''制造天花板'',即那些使复杂波前控制成为可能的几何多样性,反而成为物理器件中不确定性和性能损失的来源。

此外,超表面在动态应用中的扩展揭示了一个根本性的''速度-效率-带宽三元困境''。从相变材料、液晶到微机电系统(MEMS)执行器的主动调制机制实现了实时重构,但不可避免地会在衍射效率、光谱范围或开关速度上引入代价。目前没有任何单一平台能在三个维度上均占主导地位,迫使设计师根据应用需求接受特定的折衷。在显示和成像架构中,这种权衡被进一步加剧,因为人类感知指标(亮度、刷新率)与机器性能指标(分辨率、信噪比)之间的区别需要定制化的集成策略。应用特定部分中详细算法描述的反复出现凸显了一个系统性问题:倾向于孤立地重新优化设计参数,而不是利用统一的、基于物理的框架。

总之,超表面全息术商业化成熟度的主要障碍已不再是缺乏生成复杂设计的计算工具,而是未能将这些工具与制造约束及系统级运行要求相结合。未来的发展需要从组件级优化转向整体协同设计的范式转变。未来的进步可能依赖于联合优化元原子几何结构、制造诱导的无序性以及下游计算重建算法的框架。通过将物理可实现性和系统级指标直接嵌入逆向设计循环,研究人员可以弥合模拟最优性与实验现实之间的差距,最终实现由ODNN研究理论基础所设想的稳健、可扩展且动态的全息系统。

\bibliographystyle{unsrtnat}
\bibliography{references}
\end{document}
